\section{Implementation Details}
\label{sec:implementation_2}

The multi-task architecture uses a shared contextual representation matrix $H \in \mathbb{R}^{B \times L \times d_{\text{model}}}$ that is processed by two task-specific prediction heads. For slot filling, each token representation $H_i$ is projected onto the slot label space. For intent classification, a sentence-level representation is extracted from the backbone model and passed to a dedicated classifier.

The representation used for intent classification depends on the underlying architecture. For the custom model and BERT, the hidden state associated with the \texttt{[CLS]} token ($\mathbf{h}_{\text{[CLS]}} = H[:,0]$) is used as the sentence representation \cite{devlin2019bert}. For GPT2, which does not contain a classification token and operates autoregressively, the hidden state of the final non-padding token is selected using the attention mask ($\mathbf{h}_{\text{final}} = H[B,\sum(\text{mask})-1]$).

Training is performed using a joint objective consisting of the sum of the intent classification loss and slot filling loss:

\begin{equation}
\mathcal{L}_{\text{total}} = \mathcal{L}_{\text{intent}} + \mathcal{L}_{\text{slot}}
\end{equation}

To address the sub-tokenization introduced by WordPiece and Byte-Pair Encoding tokenizers, slot labels are assigned only to the first sub-token of each word. All remaining sub-tokens and padding positions are assigned the ignore index (\texttt{SLOT\_IGNORE\_INDEX} = -100), preventing them from contributing to the loss or gradient updates. All models are optimized using AdamW \cite{loshchilov2017decoupled} with linear learning-rate warmup and early stopping.